\documentclass[11pt]{article}
\usepackage[ngerman]{babel}
\usepackage[a4paper,margin=0.9in]{geometry}
\usepackage[T1]{fontenc}
\usepackage[utf8]{inputenc}
\usepackage{lmodern}
\usepackage{amsmath,amssymb,mathtools,mathrsfs}
\usepackage{microtype}
\usepackage{fancyhdr}
\usepackage{array,longtable,enumitem,multicol}
\usepackage{tikz}
\usetikzlibrary{patterns,calc,arrows.meta}
\setlength{\parindent}{1.5em}
\setlength{\parskip}{0.18em}
\emergencystretch=6em
\setlength{\headheight}{14pt}
\pagestyle{fancy}
\fancyhf{}
\cfoot{\thepage}
\providecommand{\sect}[2]{\section*{\S\,#1. #2}}
\providecommand{\dd}{\,d}
\providecommand{\ii}{\mathrm{i}}
\providecommand{\E}{\operatorname{E}}
\newcommand{\Z}{\mathbb Z}
\lhead{Weber Vol. I \S\S 138--147}
\rhead{Deutsche Quelle}
\begin{document}


\sect{138}{Vorzeichenbestimmung. Quadratische Reste}

Wir haben im vorigen Paragraphen zwei Formeln abgeleitet, in denen noch Vorzeichen zu bestimmen waren, nämlich
\[
\tag{1}
2^{\frac{n-1}{2}}\prod_{\nu=1}^{(n-1)/2}\cos\frac{2\nu m\pi}{n}=\pm1,
\]
\[
\tag{2}
2^{\frac{n-1}{2}}\prod_{\nu=1}^{(n-1)/2}\sin\frac{2\nu m\pi}{n}=\pm\sqrt n,
\]
worin $m$ relativ prim zu der ungeraden Zahl $n$ ist und $\nu$ die Reihe der Zahlen
\[
(\nu)\qquad 1,2,\ldots,\frac{n-1}{2}
\]
durchläuft. Wir wollen das Zeichen $\nu$ für die Zahlen dieser Reihe hier festhalten.

Jede beliebige ganze Zahl giebt bei der Theilung durch $n$ als Rest eine der Zahlen $\nu$ oder $0$ oder $n-\nu$. Statt $n-\nu$ können wir, wenn wir auch negative Reste zulassen, $-\nu$ wählen, und wenn wir also von den durch $n$ theilbaren Zahlen absehen, so bleibt eine der Zahlen $\pm\nu$ als Rest. Wir nennen diese den absolut kleinsten Rest, zum Unterschied von dem kleinsten Rest im gewöhnlichen Sinne, der aus den Zahlen $1,2,\ldots,n-1$ genommen ist.

Es sei nun also $m$ eine zu $n$ theilerfremde Zahl; wir betrachten die Reihe der Zahlen
\[
(m\nu)\qquad m,\;2m,\;3m,\ldots,\frac{n-1}{2}m
\]
und bilden zu jeder den absolut kleinsten Rest $\rho$:
\[
(\rho)\qquad \pm\nu_1,\;\pm\nu_2,\;\pm\nu_3,\ldots,
       \pm\nu_{\frac12(n-1)}.
\]
Unter diesen Zahlen $\rho$ kommen nicht zwei gleiche und auch nicht zwei entgegengesetzte vor; denn wenn die Summe oder die Differenz zweier Zahlen $(\rho)$, also $\rho+\rho'$ oder $\rho-\rho'$, gleich Null wäre, so müsste für zwei verschiedene Zahlen $\nu,\nu'$ aus $(\nu)$ die Zahl
\[
        m(\nu\pm\nu')
\]
durch $n$ theilbar sein, also müsste auch $\nu\pm\nu'$ durch $n$ theilbar sein, was unmöglich ist, da jede der Zahlen $\nu,\nu'$ kleiner als $n:2$ ist. Die Gesammtheit der Zahlen $(\rho)$ stimmt also, vom Vorzeichen und von der Reihenfolge abgesehen, mit der Gesammtheit der Zahlen $(\nu)$ überein.

In den Formeln (1) und (2) können wir, aber wegen der Periodicität von Sinus und Cosinus, $\nu m$ durch $\rho$ ersetzen, und da $\cos(-\rho)=\cos\rho$ ist, so können wir in der Formel (1) $m\nu$ auch durch $\nu$ ersetzen, d. h. das Vorzeichen ist von $m$ nicht abhängig. Um es zu bestimmen, berücksichtigen wir, dass der Cosinus eines Winkels im ersten Quadranten positiv, im zweiten Quadranten negativ ist, dass also das Vorzeichen in (1) positiv oder negativ ist, je nachdem von den Winkeln
\[
        \frac{2\nu\pi}{n}
\]
eine gerade oder eine ungerade Anzahl im zweiten Quadranten liegt, oder je nachdem eine gerade oder eine ungerade Anzahl von Zahlen $\nu$ zwischen $\frac n4$ und $\frac n2$ liegt.

Wir bezeichnen, wenn $x$ irgend eine nicht ganze Zahl ist, mit $E(x)$ die grösste ganze Zahl, die kleiner als $x$ ist, so dass $x$ zwischen $E(x)$ und $E(x)+1$ liegt. Die Anzahl der ganzen Zahlen, die zwischen zwei Zahlen $x$ und $y$ liegt, ist dann gleich $E(y)-E(x)$, und wir haben also zu untersuchen, ob
\[
        E\left(\frac n2\right)-E\left(\frac n4\right)
\]
eine gerade oder eine ungerade Zahl ist. Wir müssen vier Fälle unterscheiden:
\[
\begin{array}{lll}
 n=8k+1, & E(n/2)=4k,   & E(n/4)=2k,\\
 n=8k+3, & E(n/2)=4k+1, & E(n/4)=2k,\\
 n=8k+5, & E(n/2)=4k+2, & E(n/4)=2k+1,\\
 n=8k+7, & E(n/2)=4k+3, & E(n/4)=2k+1.
\end{array}
\]
Es ist also in (1) das positive Zeichen zu nehmen, wenn $n$ von der Form $8k+1$ oder $8k+7$, das negative, wenn $n$ von der Form $8k+3$ oder $8k+5$ ist. In den ersten Fällen ist
\[
        \frac{n^2-1}{8}
\]
eine gerade, in den beiden letzten eine ungerade Zahl, und demnach erhalten wir die genaue Formel
\[
\tag{3}
2^{\frac{n-1}{2}}\prod_{\nu=1}^{(n-1)/2}
       \cos\frac{2\nu m\pi}{n}
   =(-1)^{\frac{n^2-1}{8}}.
\]

In der Formel (2) hängt das Vorzeichen von $m$ ab; es ist positiv oder negativ, je nachdem unter den Zahlen $(\rho)$ eine gerade oder eine ungerade Anzahl negativer vorkommt. Wir setzen nun, um die Abhängigkeit des Vorzeichens von $m$ und $n$ in der Bezeichnung auszudrücken, für (2)
\[
\tag{4}
2^{\frac{n-1}{2}}\prod_{\nu=1}^{(n-1)/2}
       \sin\frac{2\nu m\pi}{n}
   =\left(\frac mn\right)\sqrt n,
\]
worin
\[
        \left(\frac mn\right)=\pm1
\]
ist. Das Symbol $\left(\frac mn\right)$ ist von einer anderen Seite her und in speciellerer Fassung von Legendre in die Zahlentheorie eingeführt und von Jacobi verallgemeinert worden. Wir wollen es das Legendre'sche Symbol nennen. Seine Bedeutung für die Zahlentheorie wird sich gleich ergeben; zunächst haben wir den Satz:

\begin{enumerate}[label=\arabic*.,leftmargin=2em]
\item Es ist $\left(\frac mn\right)$ gleich $+1$ oder gleich $-1$, je nachdem die Anzahl der negativen unter den absolut kleinsten Resten von $m\nu$ eine gerade oder eine ungerade ist.

Für $m=1$ gilt, wie wir schon gesehen haben, das positive Zeichen in der Formel (2); also haben wir, wie auch die Formel (4) direct zeigt,
\[
        \left(\frac 1n\right)=+1.
\]

\item Verwandeln wir $m$ in $-m$, so ändert sich in allen Factoren des Productes auf der linken Seite von (4) das Vorzeichen; da die Anzahl der Factoren $\frac{n-1}{2}$ beträgt, so ergiebt sich
\[
\tag{3'}
        \left(\frac{-m}{n}\right)=(-1)^{\frac{n-1}{2}}\left(\frac mn\right),
\]
und für $m=1$
\[
\tag{4'}
        \left(\frac{-1}{n}\right)=(-1)^{\frac{n-1}{2}}.
\]

\item Setzen wir $2m$ für $m$ und benutzen die Formel
\[
        \sin\frac{4\nu m\pi}{n}
        =2\sin\frac{2\nu m\pi}{n}\cos\frac{2\nu m\pi}{n},
\]
so folgt aus (3)
\[
\tag{5}
        \left(\frac{2m}{n}\right)=(-1)^{\frac{n^2-1}{8}}\left(\frac mn\right),
\]
und für $m=1$
\[
\tag{6}
        \left(\frac 2n\right)=(-1)^{\frac{n^2-1}{8}}.
\]

\item Aendert man $m$ um ein Vielfaches von $n$, so bleiben alle Factoren des Productes (4) ungeändert, und daraus folgt
\[
\tag{7}
        \left(\frac mn\right)=\left(\frac {m'}n\right),
        \qquad\text{wenn }m\equiv m'\pmod n.
\]
\end{enumerate}

Wir setzen jetzt nicht nur $n$, sondern auch $m$ als ungerade und positiv voraus und wenden die Formel (29) des §.\,137 an, in der wir $n$ durch $m$ ersetzen:
\[
\tag{5}
\frac{\sin m\varphi}{\sin\varphi}
=(-1)^{\frac{m-1}{2}}2^{m-1}\prod_\beta(\sin^2\varphi-\beta),
\]
worin $\beta$ die sämmtlichen Wurzeln der Gleichung $\Phi_m(x)=0$ durchläuft, die den Ausdruck
\[
\tag{6}
        \beta=\sin^2\frac{2\mu\pi}{m}
\]
haben. Ist nun $m'$ eine zu $n$ theilerfremde Zahl, so durchläuft
\[
\tag{7}
        \alpha=\sin^2\frac{2\nu m'\pi}{n}
\]
die sämmtlichen Wurzeln der Gleichung $\Phi_n(x)=0$, und wenn wir in (5)
\[
\tag{8}
        \varphi=\frac{2\nu m'\pi}{n}
\]
setzen und das Product in Bezug auf alle $\nu$ nehmen, so folgt
\[
\tag{9}
\frac{\prod_\nu \sin\frac{2\nu m m'\pi}{n}}
     {\prod_\nu \sin\frac{2\nu m'\pi}{n}}
=(-1)^{\frac{(m-1)(n-1)}4}
2^{\frac{(m-1)(n-1)}2}\prod_\alpha\prod_\beta(\alpha-\beta).
\]
Die rechte Seite dieser Formel ist aber von $m'$ ganz unabhängig. Der Werth der linken Seite ist also derselbe, als ob $m'=1$ wäre, und wenn man für die Producte die Werthe aus (4) setzt, so folgt
\[
        \left(\frac{mm'}n\right):\left(\frac{m'}n\right)
        =
        \left(\frac mn\right):\left(\frac 1n\right),
\]
oder
\[
        \left(\frac{mm'}n\right)=
        \left(\frac mn\right)\left(\frac{m'}n\right).
\]
Hier war zunächst vorausgesetzt, dass $m$ und $m'$ positiv und ungerade seien. Nach 3. und 5. aber bleibt diese Formel auch noch richtig, wenn $m$ oder $m'$ negativ oder gerade ist, wenn nur $m$ und $m'$ relativ prim zu $n$ sind.

Nehmen wir wieder $m$ ungerade und positiv an, so folgt jetzt aus (4) und (9)
\[
\tag{10}
\left(\frac mn\right)
=(-1)^{\frac{(m-1)(n-1)}4}
2^{\frac{(m-1)(n-1)}2}\prod_\alpha\prod_\beta(\alpha-\beta).
\]
Wenn wir hierin $m$ mit $n$ vertauschen, so ändert in dem Doppelproduct auf der rechten Seite jeder Factor sein Vorzeichen; alles Andere bleibt ungeändert. Die Anzahl dieser Factoren ist aber
\[
        \frac{n-1}{2}\cdot\frac{m-1}{2},
\]
und daraus folgt
\[
\tag{9'}
        \left(\frac mn\right)
        =(-1)^{\frac{(m-1)(n-1)}4}\left(\frac nm\right).
\]
Oder: $\left(\frac mn\right)$ ist gleich $\left(\frac nm\right)$, wenn von den beiden ungeraden Zahlen $m,n$ wenigstens eine von der Form $4k+1$ ist, und $\left(\frac mn\right)$ ist entgegengesetzt zu $\left(\frac nm\right)$, wenn beide Zahlen von der Form $4k+3$ sind.

Dieser berühmte Satz ist unter dem Namen des Reciprocitätsgesetzes bekannt. Der hier gegebene Beweis rührt von Eisenstein her. Man kann mit seiner Hülfe und nach 5. und 6. den Werth des Symbols $\left(\frac mn\right)$ sehr schnell ermitteln, indem man so verfährt, als ob es sich um die Bestimmung des grössten gemeinschaftlichen Theilers von $m$ und $n$ handelte.

Aus 8. und 9. ergiebt sich noch ein letzter Satz, der gilt, wenn $n$ und $n'$ zwei positive ungerade zu $m$ theilerfremde Zahlen sind:
\[
\tag{10'}
        \left(\frac mn\right)\left(\frac m{n'}\right)
        =
        \left(\frac m{nn'}\right).
\]
Seine Richtigkeit ergiebt sich, wenn $m$ ungerade und positiv ist, auf Grund der Congruenz
\[
\tag{7'}
 (nn'-1)-(n-1)-(n'-1)
 = (n-1)(n'-1)\equiv0\pmod 4,
\]
wenn man auf beide Seiten von 10. die Formel 9. und dann 8. anwendet, und wenn $m$ gerade oder negativ ist, aus 4. und 6.

Nach diesen Sätzen kann man $\left(\frac mn\right)$ auf ein Product von Werthen $\left(\frac qp\right)$ zurückführen, worin $p,q$ Primzahlen sind. Die Berechnung von $\left(\frac mn\right)$ geschieht aber leichter nach 9. ohne diese Zerlegung.

Wir wollen nun noch eine andere Bedeutung des Symbols $\left(\frac mp\right)$ kennen lernen, für den Fall, dass $p$ eine ungerade Primzahl ist, durch die ursprünglich das Symbol von Legendre eingeführt ist.

Wenn wir alle durch $p$ nicht theilbaren ganzen Zahlen $x$ ins Quadrat erheben und die Reste der Division durch $p$ aufsuchen, so erhalten wir im Ganzen nur $\frac{p-1}{2}$ verschiedene Reste; denn erstens geben die Quadrate congruenter Zahlen dieselben Reste, und zweitens geben die Quadrate entgegengesetzter Zahlen auch dieselben Reste. Wir bekommen also gewiss alle Reste, wenn wir die $\frac{p-1}{2}$ Zahlen
\[
(\nu)\qquad 1,2,3,\ldots,\frac{p-1}{2}
\]
ins Quadrat erheben. Auf der anderen Seite ergeben auch die Quadrate dieser Zahlen $\nu$ lauter verschiedene Reste; denn es kann, wenn $\nu$ und $\nu'$ verschiedene von ihnen sind, niemals
\[
        \nu^2-\nu'^2=(\nu-\nu')(\nu+\nu')
\]
durch $p$ theilbar sein, weil sowohl $\nu-\nu'$ als $\nu+\nu'$ kleiner als $p$ ist. Unter den Zahlen
\[
(s)\qquad 1,2,3,\ldots,p-1
\]
kommt also die Hälfte unter den Resten von $x^2$ vor, die Hälfte nicht. Die ersteren Zahlen und alle mit ihnen nach $p$ congruenten heissen die quadratischen Reste von $p$, die anderen die quadratischen Nichtreste von $p$.

Wenn nun $m$ quadratischer Rest ist, so ist die Congruenz
\[
        x^2\equiv m\pmod p
\]
möglich, und aus 7. und 8. folgt
\[
        \left(\frac mp\right)
        =\left(\frac{x^2}p\right)
        =\left(\frac xp\right)^2=+1.
\]
Es ist also $\left(\frac mp\right)=+1$, wenn $m$ quadratischer Rest von $p$ ist. Es ist aber in §.\,135, 1. der Satz ausgesprochen, dass, wenn für $m$ die Reihe der Zahlen $1,2,\ldots,p-1$ gesetzt wird, $\left(\frac mp\right)$ ebenso oft das positive wie das negative Zeichen hat, und daraus ersieht man, dass das negative Zeichen gilt, wenn $m$ quadratischer Nichtrest ist, es folgt also:

\begin{enumerate}[label=\arabic*.,leftmargin=2em,start=11]
\item Ist $p$ eine ungerade Primzahl und $m$ durch $p$ nicht theilbar, so ist
\[
        \left(\frac mp\right)=+1\quad\hbox{oder}\quad -1,
\]
je nachdem $m$ quadratischer Rest oder quadratischer Nichtrest von $p$ ist.
\end{enumerate}

Darin sind alle Hauptsätze über die quadratischen Reste enthalten, z. B. der Satz: Das Product zweier quadratischer Reste oder zweier Nichtreste ist ein Rest; das Product aus einem Rest und einem Nichtrest ist ein Nichtrest.

\section*{Drittes Buch.\\Algebraische Grössen}

\section*{Dreizehnter Abschnitt.\\Die Galois'sche Theorie}

\sect{139}{Der Körperbegriff}

Ein System von Zahlen wird ein Zahlkörper genannt, wenn es so in sich vollendet und abgeschlossen ist, dass die vier fundamentalen Rechenoperationen, die Addition, die Subtraction, die Multiplication und die Division, ausgeführt mit irgend welchen Zahlen des Systems, ausgenommen die Division durch Null, immer auf Zahlen führen, die demselben System angehören. Dieser Begriff, der eine Eintheilung der Zahlenarten nach einem natürlichen Gesichtspunkte giebt, ist von Dedekind eingeführt. Er ist für die Algebra von der grössten Bedeutung, und es ist nicht gleichgültig, dafür einen bezeichnenden und ausdrucksvollen Namen zu haben. Das Wort Zahlkörper ist von Dedekind nach zahlreichen Analogien gebildet, in denen das Wort Körper in ähnlicher Weise eine Vereinigung von zusammengehörigen Dingen bezeichnet, der eine gewisse Vollständigkeit zukommt.

Der Begriff des Zahlkörpers kann erweitert und auf alle Grössen übertragen werden, mit denen nach den Regeln der vier Species gerechnet werden kann; insbesondere also auf rationale Functionen irgend welcher Veränderlichen.

Wir nennen also einen Functionenkörper ein System von Functionen von einer oder mehreren Veränderlichen von der Beschaffenheit, dass in diesem System die Rechnungen mit den vier Species unbegrenzt ausgeführt werden können und immer auf eine bestimmte Function desselben Systems führen, immer mit Ausnahme der Division durch Null. Die Veränderlichen können ganz von einander unabhängig sein; es ist aber auch der Fall nicht ausgeschlossen, dass gewisse Beziehungen zwischen ihnen festgesetzt sind, die beim Rechnen zu berücksichtigen sind.

Eine Function eines Functionenkörpers gilt nur dann als Null, wenn sie identisch, d. h. für alle in Betracht kommenden Werthe der Veränderlichen, verschwindet.

Da wir vorläufig unsere Betrachtungen nicht einschränken wollen, so werden wir jetzt von Körpern schlechtweg sprechen, und die Objecte, mit denen die Rechnungen auszuführen sind, die sowohl Zahlen als Functionen sein können, als Grössen oder auch als die Elemente des Körpers bezeichnen.

Ein Körper ist dann also ein System von Grössen von der Vollständigkeit, dass in ihm die Grössen addirt, subtrahirt, multiplicirt und dividirt werden können. Wenn alle Grössen eines Körpers in einem zweiten Körper enthalten sind, so heisst der erste Körper ein Theiler des zweiten. Ist $\Omega$ ein Theiler von $\Omega'$, so werden wir auch sagen, $\Omega'$ ist ein Körper über $\Omega$.

Streng genommen, bildet die einzige Zahl Null für sich einen Körper. Diesen wollen wir aber der Kürze wegen ein für allemal ausnehmen. Dann ist das nächstliegende Beispiel eines Körpers der Inbegriff aller rationalen Zahlen. Dieser Körper ist ein Theiler von jedem anderen Körper; denn jeder Körper enthält wenigstens eine von Null verschiedene Grösse $\omega$, also auch den Quotienten $\omega:\omega=1$, d. h. die Zahl $1$, also auch alle ganzen Zahlen; und aus den ganzen Zahlen kann man wieder durch Division alle Brüche ableiten.

Andere Beispiele von Zahlkörpern sind: der Inbegriff aller complexen Zahlen von der Form $x+yi$, worin $i=\sqrt{-1}$, $x$ und $y$ alle rationalen Zahlen bedeuten; ferner der Inbegriff aller reellen Zahlen, oder der Inbegriff aller überhaupt existirenden complexen Zahlen $x+yi$.

Als Beispiel eines Functionenkörpers mag der Inbegriff aller ganzen und gebrochenen rationalen Functionen einer Veränderlichen, mit Einschluss der Constanten, dienen, wobei man die Constanten auf einen beliebigen Zahlkörper beschränken oder auch ganz unbeschränkt annehmen kann.

\sect{140}{Adjunction}

Wenn einem Körper irgend welcher Grössen, den wir mit $\Omega$ bezeichnen wollen, irgend eine Grösse $\alpha$ hinzugefügt wird, die nicht in ihm enthalten ist, so entsteht ein neues Grössensystem
\[
        \Omega,\alpha,
\]
das aber kein Körper sein wird; um daraus einen Körper abzuleiten, muss man alle Grössen hinzufügen, die sich durch die Verbindung von $\alpha$ mit den Zahlen von $\Omega$ mittelst der Grundrechnungsarten ableiten lassen. So erhält man einen erweiterten Körper $\Omega'$, der zugleich $\alpha$ und $\Omega$ enthält, und der durch $\alpha$ und $\Omega$ völlig bestimmt ist. Wir wollen ihn den Körper $\Omega,\alpha$ nennen.

Dies Hinzufügen einer neuen Grösse $\alpha$ zu einem Körper $\Omega$ heisst Adjungiren, und man sagt, der Körper $\Omega'$ entsteht aus $\Omega$ durch Adjunction von $\alpha$.

Zu $\Omega'$ kann man wieder eine Grösse $\alpha'$ adjungiren und erhält einen dritten Körper $\Omega''$, der dann $\Omega,\alpha,\alpha'$ enthält. Man kann aber gleichzeitig zwei oder mehrere Zahlen zu $\Omega$ adjungiren und erhält so denselben Körper, wenn man gleichzeitig $\alpha$ und $\alpha'$ dem $\Omega$ adjungirt. Man kann auch einem Körper einen zweiten Körper adjungiren und erhält so einen Körper, der zugleich die Zahlen der beiden Körper enthält. Adjungirt man einem Körper einen seiner Theile, so entsteht kein neuer Körper.

So erhält man z. B., wenn man zum Körper der rationalen Zahlen, der mit $\mathfrak R$ bezeichnet sein mag, die Zahl $i=\sqrt{-1}$ adjungirt, den Körper der complexen Zahlen $x+yi$ mit rationalen $x,y$, den wir $\mathfrak C$ nennen wollen. Durch Adjunction einer Variablen $u$ zum Körper $\mathfrak R$ erhält man den Körper der rationalen Functionen von $u$ mit rationalen Coefficienten, durch Adjunction von $u$ und $i$ zu $\mathfrak R$ den Körper der rationalen Functionen von $u$, deren Coefficienten Zahlen in $\mathfrak C$ sind, u. s. f.

Es sei schliesslich noch bemerkt, dass ein Körper von Grössen vielfach auch, nach Kronecker's Vorgang, ein Rationalitätsbereich genannt wird. Es ist bisweilen nützlich, den Ausdruck Rationalitätsbereich für einen Körper dann zu brauchen, wenn damit ausgedrückt werden soll, dass bei einer bestimmten Aufgabe die Grössen dieses Körpers als bekannt oder rational betrachtet werden sollen.

\sect{141}{Functionen in einem Körper}

Wir haben oft Veranlassung, ganze rationale Functionen einer Veränderlichen $x$
\[
\tag{1}
        f(x)=a_0x^m+a_1x^{m-1}+\cdots+a_{m-1}x+a_m
\]
zu betrachten, deren Coefficienten einem bestimmten Körper $\Omega$ angehören, und die wir in $\Omega$ enthaltene Functionen, oder kurz Functionen in $\Omega$ nennen wollen. Sie sind nicht zu verwechseln mit Functionen, die etwa dem Körper $\Omega$ angehören, wenn $\Omega$ ein Functionenkörper ist.

Im dritten Abschnitt haben wir gesehen, dass, wenn die Coefficienten bestimmte Zahlenwerthe haben, $m$ Wurzeln $\alpha_1,\alpha_2,\ldots,\alpha_m$ von $f(x)$ existiren, die gleichfalls bestimmte Zahlenwerthe haben. Sind die Coefficienten unabhängig veränderlich, so giebt es für jedes Werthsystem derselben ein Werthsystem der Wurzeln, und die Wurzeln ändern sich stetig mit den Coefficienten, so lange man die Veränderung auf ein hinlänglich enges Gebiet beschränkt, in dem die Discriminante nicht verschwindet.

Wenn nun $\Omega$ ein beliebiger Körper und $f(x)$ eine Function in $\Omega$ ist, so sind zwei Fälle möglich. Die Function $f(x)$ ist entweder zerlegbar oder nicht zerlegbar in Factoren niedrigeren Grades, deren Coefficienten dem Körper $\Omega$ angehören. Im ersten Falle heisst $f(x)$ reducibel, im zweiten irreducibel in $\Omega$. Ist der Körper $\Omega$ als bekannt und feststehend betrachtet, so spricht man bisweilen schlechtweg von reducibeln oder irreducibeln Functionen.

Eine lineare Function ist selbstverständlich in jedem Körper irreducibel. Multiplicirt man mehrere Functionen der Form (1) mit einander, so entsteht eine Function derselben Form, die reducibel ist. Eine Function, die in $\Omega$ irreducibel ist, kann in einem erweiterten Körper reducibel werden; so wird jede Function $f(x)$ reducibel in dem Körper, der aus $\Omega$ durch Adjunction einer Wurzel $\alpha$ von $f(x)=0$ entsteht.

\begin{enumerate}[label=\Roman*.,leftmargin=2.8em]
\item Eine irreducible Function $f(x)$ kann mit einer anderen Function $F(x)$, deren Coefficienten demselben Körper angehören, keinen gemeinsamen Theiler haben, wenn nicht $F(x)$ durch $f(x)$ theilbar ist.

Dieser Satz ist fast selbstverständlich; denn der grösste gemeinschaftliche Theiler zweier Functionen $F(x)$ und $f(x)$ lässt sich durch rationale Rechnung finden und ist daher auch in $\Omega$ enthalten. Da nun aber $f(x)$ keinen in $\Omega$ enthaltenen Theiler hat als sich selbst oder eine Constante, so muss dieser grösste gemeinschaftliche Theiler entweder eine Constante oder $f(x)$ selbst sein.

Es folgt daraus, dass eine irreducible Function niemals mehrfache Factoren haben kann; denn sie müsste sonst mit ihrer Ableitung $f'(x)$ einen gemeinsamen Theiler haben.

\item Ist $f(x)$ irreducibel, und verschwindet $F(x)$ für eine Wurzel von $f(x)=0$, so verschwindet es auch für alle anderen Wurzeln von $f(x)=0$.

\item Ist $F(x)$ von niedrigerem Grade als die irreducible Function $f(x)$, und verschwindet $F(x)$ für eine Wurzel von $f(x)$, so muss $F(x)$ identisch verschwinden.
\end{enumerate}

Auch ganze rationale Functionen von mehreren Variablen $f(x,y,z,\ldots)$, deren Coefficienten dem Körper $\Omega$ angehören, heissen Functionen in $\Omega$. Auch diese werden als reducible und irreducible Functionen bezeichnet, je nachdem sie zerlegbar oder nicht zerlegbar sind in mehrere Factoren, die selbst Functionen in $\Omega$ sind und wenigstens einige der Variablen wirklich enthalten. Hier ist aber auf einen Unterschied aufmerksam zu machen: es giebt unzerlegbare Functionen, ferner zerlegbare Functionen, deren Factoren nicht in $\Omega$ liegen, und endlich Functionen, die in mehrere Functionen in $\Omega$ zerlegbar sind; nur die letztere Art werden wir reducibel in $\Omega$ nennen.

Die fundamentalen Sätze über irreducible Functionen haben wir in §.\,51 kennen gelernt; benutzt ist dort nur die Reducibilität oder Irreducibilität einer Function, und die besondere Natur des Körpers $\Omega$ kommt nicht in Betracht. So ist darin der Beweis des Satzes enthalten:

\begin{enumerate}[label=\Roman*.,leftmargin=2.8em,start=4]
\item Eine reducible Function in $\Omega$ kann nur auf eine Art in irreducible Factoren zerlegt werden, wenn man rationale Functionen, die sich nur durch einen constanten Factor unterscheiden, als nicht verschieden betrachtet.
\end{enumerate}

\sect{142}{Algebraische Körper}

Ist $\Omega$ ein beliebiger Körper, und $F(x)$ eine Function in $\Omega$, so heisst die Gleichung
\[
\tag{1}
        F(x)=0
\]
eine Gleichung in $\Omega$. Diese Gleichung heisst reducibel oder irreducibel, je nachdem die Function $F(x)$ reducibel oder irreducibel ist. Durch Adjunction einer Wurzel $\alpha$ einer solchen Gleichung zu $\Omega$ entsteht, wenn $\alpha$ nicht selbst schon zu $\Omega$ gehört, ein neuer Körper $\Omega'$, den wir einen algebraischen Körper über $\Omega$ nennen wollen. Wir brauchen für einen solchen Körper das Zeichen
\[
        \Omega(\alpha).
\]

Sind $\beta,\gamma,\ldots$ Wurzeln von anderen Gleichungen in $\Omega$, oder auch andere Wurzeln derselben Gleichung, so erhalten wir durch gleichzeitige Adjunction von $\alpha,\beta,\gamma,\ldots$ gleichfalls algebraische Körper über $\Omega$, die wir mit
\[
        \Omega(\alpha,\beta,\gamma,\ldots)
\]
bezeichnen. Wir werden gleich sehen, dass diese Erweiterung des Begriffes algebraischer Körper nur eine scheinbare ist.

Die Gleichung (1), deren Wurzel $\alpha$ ist, kann reducibel sein. Unter den irreducibeln Factoren von $F(x)$ ist aber wenigstens einer, der für $x=\alpha$ verschwindet, den wir mit $f(x)$ bezeichnen wollen. Diese Function $f(x)$ ist, wenn der Coefficient der höchsten Potenz von $x$ gleich $1$ angenommen wird, völlig bestimmt. Die Gleichung $f(x)=0$ hat also die Form
\[
\tag{2}
        x^n+a_1x^{n-1}+a_2x^{n-2}+\cdots+a_{n-1}x+a_n=0,
\]
worin $a_1,a_2,\ldots,a_n$ Grössen in $\Omega$ sind. Ist $n$ der Grad dieser Gleichung, so nennen wir $n$ auch den Grad des Körpers $\Omega(\alpha)$.

Alle Grössen $\Theta$ des Körpers $\Omega(\alpha)$ lassen sich ableiten durch Addition, Subtraction, Multiplication und Division aus $\alpha$ und aus Zahlen in $\Omega$; sie lassen sich also darstellen als rationale Functionen von $\alpha$ mit Coefficienten in $\Omega$, oder in der Form
\[
\tag{3}
        \Theta=\frac{\chi(\alpha)}{\psi(\alpha)},
\]
worin $\chi(x)$ und $\psi(x)$ Functionen in $\Omega$ sind. Da aber $\psi(\alpha)$ von Null verschieden sein muss, so ist $\psi(x)$ nicht durch $f(x)$ theilbar, und da $f(x)$ irreducibel angenommen ist, so ist $\psi(x)$ relativ prim zu $f(x)$.

Danach können wir nach §.\,6, Satz II die Functionen $Q(x)$, $\varphi(x)$ in $\Omega$, und zwar $\varphi(x)$ höchstens vom Grade $n-1$, so bestimmen, dass
\[
\tag{4}
        Q(x)f(x)+\varphi(x)\psi(x)=\chi(x)
\]
wird; setzen wir hierin $x=\alpha$, so dass $f(\alpha)=0$ wird, so folgt
\[
\tag{5}
        \frac{\chi(\alpha)}{\psi(\alpha)}=\varphi(\alpha).
\]
Es gilt also:

\begin{enumerate}[label=\arabic*.,leftmargin=2em]
\item Jede Grösse des algebraischen Körpers $\Omega(\alpha)$ kann auf eine und nur eine Weise in der Form
\[
        \Theta=c_0+c_1\alpha+c_2\alpha^2+\cdots+c_{n-1}\alpha^{n-1}
\]
dargestellt werden, worin die Coefficienten $c_0,c_1,\ldots,c_{n-1}$ Grössen in $\Omega$ sind.
\end{enumerate}

Dass diese Darstellung nur auf eine Art möglich ist, folgt aus Satz III des vorigen Paragraphen: wenn zwei Darstellungen vorhanden wären, so ergäbe ihre Differenz eine Function niedrigeren Grades als $f(x)$, die für $x=\alpha$ verschwindet, also identisch verschwinden müsste.

Wenn wir die Grössen des Körpers $\Omega$ als die gegebenen, seine Grössen als rational, den Körper $\Omega(\alpha)$ als den daraus abgeleiteten betrachten, so bezeichnen wir eine rationale Function in $\Omega$ auch kurz als eine rationale Function, und können demnach auch sagen:
\[
        \hbox{Der Körper }\Omega(\alpha)\hbox{ besteht aus allen rationalen Functionen von }\alpha.
\]

\sect{143}{Mehrfache Adjunction}

Einer genaueren Untersuchung des Körpers $\Omega(\alpha)$ schicken wir einen allgemeinen, sehr einfachen, aber folgereichen Satz voraus.

\begin{enumerate}[label=\arabic*.,leftmargin=2em]
\item Sind $\Phi_1(x,y,z,\ldots),\Phi_2(x,y,z,\ldots),\Phi_3(x,y,z,\ldots),\ldots$ ganze rationale Functionen der Veränderlichen $x,y,z,\ldots$ mit beliebigen Coefficienten, die in keiner der Functionen alle zugleich verschwinden, so kann man für die Veränderlichen auf unendlich viele Arten solche rationale Zahlwerthe setzen, dass keine der Functionen $\Phi_1,\Phi_2,\Phi_3,\ldots$ verschwindet.
\end{enumerate}

Der Satz ist zunächst evident, wenn die Functionen nur von einer Veränderlichen abhängen. Dann aber kann man aus dem Satz für $n$ Veränderliche den Satz für $n+1$ Veränderliche herleiten, indem man die Functionen nach der $(n+1)$ten Veränderlichen ordnet, die übrigen $n$ Veränderlichen so wählt, dass in keiner Function alle Coefficienten verschwinden, und dann für die eine übrig bleibende Veränderliche einen geeigneten rationalen Werth wählt.

Es soll jetzt zunächst nachgewiesen werden, dass man die gleichzeitige Adjunction mehrerer algebraischen Grössen durch die Adjunction einer einzigen ersetzen kann. Es seien also $\alpha,\beta,\gamma,\ldots$ algebraische Grössen und
\[
\tag{1}
        A(x)=0,\qquad B(x)=0,\qquad C(x)=0,\ldots
\]
Gleichungen in $\Omega$, deren Wurzeln $x=\alpha,x=\beta,x=\gamma,\ldots$ sind. Keine dieser Gleichungen soll eine mehrfache Wurzel haben.

Jede Grösse des Körpers $\Omega(\alpha,\beta,\gamma,\ldots)$ kann als ganze rationale Function in $\Omega$ von $\alpha,\beta,\gamma,\ldots$ dargestellt werden.

Wir bilden eine lineare Function der $\alpha,\beta,\gamma,\ldots$
\[
        \xi=x\alpha+y\beta+z\gamma+\cdots
\]
und beachten, dass jede der Gleichungen (1) mehrere Wurzeln hat. Ist also $\alpha',\beta',\gamma',\ldots$ irgend eine von $\alpha,\beta,\gamma,\ldots$ verschiedene Combination von je einer Wurzel von jeder der Gleichungen (1), so setzen wir
\[
        \xi'=x\alpha'+y\beta'+z\gamma'+\cdots,
\]
und bilden ebenso $\xi'',\xi''',\ldots$. Die Anzahl der so gebildeten Grössen $\xi$ ist, wenn $a,b,c,\ldots$ die Grade von $A(x),B(x),C(x),\ldots$ sind,
\[
        m=abc\cdots.
\]
Die Differenzen $\xi-\xi'$, $\xi-\xi''$, $\xi'-\xi''$, $\ldots$ sind lineare Functionen von $x,y,z,\ldots$, und keine ist identisch Null. Nach dem Satz 1 können wir also für $x,y,z,\ldots$ rationale Zahlen setzen, dass die $m$ Werthe $\xi,\xi',\xi'',\ldots$ alle von einander verschieden sind.

Nun ist jede rationale Function, die in Bezug auf die Wurzeln jeder der Gleichungen (1) symmetrisch ist, rational durch die Coefficienten dieser Gleichungen ausdrückbar, also eine Zahl in $\Omega$. Dazu gehören auch die Coefficienten der Function $m$ten Grades von $t$:
\[
\tag{2}
        F(t)=(t-\xi)(t-\xi')(t-\xi'')\cdots,
\]
und $F(t)=0$ ist also eine Gleichung in $\Omega$, die keine gleichen Wurzeln hat, und deren eine Wurzel $\xi$ ist.

Ist ferner $\Theta$ eine Grösse in $\Omega(\alpha,\beta,\gamma,\ldots)$, also eine ganze rationale Function von $\alpha,\beta,\gamma,\ldots$, und bezeichnen wir mit $\Theta',\Theta'',\ldots$ die Grössen, die aus dieser rationalen Function hervorgehen, wenn die Variablen durch $\alpha',\beta',\gamma',\ldots$ oder $\alpha'',\beta'',\gamma'',\ldots$ ersetzt werden, so ist
\[
        F(t)\left(
        \frac{\Theta}{t-\xi}+
        \frac{\Theta'}{t-\xi'}+
        \frac{\Theta''}{t-\xi''}+\cdots
        \right)=\psi(t)
\]
als symmetrische Function der Wurzeln eine ganze rationale Function $(m-1)$ten Grades von $t$ in $\Omega$, und wenn man $t=\xi$ setzt, so folgt
\[
\tag{3}
        \Theta=\frac{\psi(\xi)}{F'(\xi)}.
\]
Also ist $\Theta$ in $\Omega(\xi)$ enthalten. Da umgekehrt jede Grösse in $\Omega(\xi)$ zugleich in $\Omega(\alpha,\beta,\gamma,\ldots)$ enthalten ist, so sind beide Körper identisch:
\[
        \Omega(\xi)=\Omega(\alpha,\beta,\gamma,\ldots).
\]
Damit ist unsere Behauptung erwiesen.

\sect{144}{Primitive und imprimitive Körper}

Nach dem zuletzt bewiesenen Satze beschränken wir uns jetzt auf die Betrachtung algebraischer Körper $\Omega(\alpha)$, worin $\alpha$ die Wurzel einer Gleichung in $\Omega$ ist. Diese Gleichung möge, von mehrfachen Factoren befreit, mit
\[
\tag{1}
        F(x)=0
\]
bezeichnet und vom Grade $m$ sein. Ihre Wurzeln seien
\[
\tag{2}
        \alpha,\alpha_1,\alpha_2,\ldots,\alpha_{m-1}.
\]
Es kann $F(x)$ in $\Omega$ reducibel sein; dann wird es einen und nur einen irreducibeln Factor $f(x)$ von $F(x)$ geben, so dass $\alpha$ eine Wurzel der Gleichung
\[
\tag{3}
        f(x)=0
\]
ist. Der Grad von $f(x)$ sei $n$, und die Wurzeln von (3), die alle unter den Wurzeln (2) enthalten sind, seien
\[
\tag{4}
        \alpha,\alpha_1,\alpha_2,\ldots,\alpha_{n-1}.
\]
Der Grad des Körpers $\Omega(\alpha)$ ist dann $n$. Jede der Grössen (4) giebt zu einem algebraischen Körper Anlass, und so entstehen die Körper
\[
\tag{5}
        \Omega(\alpha),\Omega(\alpha_1),\ldots,\Omega(\alpha_{n-1}),
\]
die wir die mit $\Omega(\alpha)$ conjugirten Körper nennen.

Nach §.\,142 erhalten wir jede Grösse $\Theta$ in $\Omega(\alpha)$ in der Form
\[
        \Theta=\chi(\alpha),
\]
wobei $\chi(t)$ eine ganze rationale Function in $\Omega$ höchstens vom Grade $n-1$ ist. Setzen wir für $t$ die verschiedenen Grössen (4), so erhalten wir
\[
\tag{6}
        \Theta=\chi(\alpha),\quad
        \Theta_1=\chi(\alpha_1),\quad\ldots,\quad
        \Theta_{n-1}=\chi(\alpha_{n-1}),
\]
und diese heissen die mit $\Theta$ conjugirten Grössen. Jede symmetrische Function dieser Grössen ist zugleich eine symmetrische Function der Wurzeln der Gleichung (3) und mithin in $\Omega$ enthalten.

Unter diesen symmetrischen Functionen wollen wir zwei hervorheben: die Summe
\[
\tag{7}
        S(\Theta)=\Theta+\Theta_1+\Theta_2+\cdots+\Theta_{n-1},
\]
die wir die Spur von $\Theta$ nennen, und das Product
\[
\tag{8}
        N(\Theta)=\Theta\Theta_1\Theta_2\cdots\Theta_{n-1},
\]
das die Norm von $\Theta$ heisst. Conjugirte Zahlen haben hiernach dieselben Spuren und Normen.

Das Product
\[
\tag{9}
        (t-\Theta)(t-\Theta_1)\cdots(t-\Theta_{n-1})=\Phi(t)
\]
ist eine ganze rationale Function $n$ten Grades von $t$ in $\Omega$, und ihre Wurzeln sind $\Theta$ und die mit $\Theta$ conjugirten Grössen. Daraus ergiebt sich:
\begin{enumerate}[label=\arabic*.,leftmargin=2em]
\item Jede Grösse in $\Omega(\alpha)$ ist Wurzel einer Gleichung $n$ten Grades in $\Omega$, deren übrige Wurzeln die mit $\Theta$ conjugirten Grössen sind.
\end{enumerate}

Die Berechnung der Function $\Phi(t)$ aus $f(x)$ ist nichts Anderes als die Tschirnhausen-Transformation.

Wir haben nun die Function $\Phi(t)$ auf ihre Irreducibilität zu untersuchen. Wenn $\Phi(t)$ reducibel ist, so hat sie einen irreducibeln Factor $\varphi(t)$ in $\Omega$. Ist
\[
        \varphi(\Theta)=\varphi[\chi(\alpha)]=0,
\]
so haben die Gleichungen $\varphi[\chi(x)]=0$ und $f(x)=0$ eine gemeinsame Wurzel. Da aber $f(x)$ irreducibel ist, so muss $\varphi[\chi(x)]$ für alle Wurzeln $\alpha,\alpha_1,\ldots,\alpha_{n-1}$ verschwinden. Sind die mit $\Theta$ conjugirten Werthe alle von einander verschieden, so ist $\varphi(t)$ mit $\Phi(t)$ identisch; andernfalls seien unter den conjugirten Werthen nur $n_1$ verschiedene vorhanden. Dann ist
\[
\tag{10}
        \varphi(t)=(t-\Theta)(t-\Theta_1)\cdots(t-\Theta_{n_1-1}),
\]
und
\[
\tag{11}
        \Phi(t)=\varphi(t)^{n_2},
\]
\[
\tag{12}
        n=n_1n_2.
\]
Daraus ergiebt sich der Satz:
\begin{enumerate}[label=\arabic*.,leftmargin=2em,start=2]
\item Die Function $\Phi(t)$ ist entweder irreducibel oder sie ist eine Potenz einer irreducibeln Function. Die $n$ mit einer Zahl in $\Omega(\alpha)$ conjugirten Zahlen sind entweder alle von einander verschieden, oder sie zerfallen in $n_1$ Systeme von je $n_2$ unter einander gleichen Zahlen. Im ersten Falle ist $\Phi(t)$ irreducibel, im zweiten die $n_2$te Potenz einer irreducibeln Function $n_1$ten Grades in $\Omega$.
\end{enumerate}

Eine Grösse $\Theta$ in $\Omega(\alpha)$, die von allen ihren conjugirten Zahlen verschieden ist und die also einer irreducibeln Gleichung $n$ten Grades genügt, heisst eine primitive Grösse des Körpers. Nach §.\,143 lassen sich unendlich viele solche primitive Grössen bestimmen.

\begin{enumerate}[label=\arabic*.,leftmargin=2em,start=3]
\item Jede Grösse $\omega$ des Körpers $\Omega(\alpha)$ kann rational durch eine beliebige primitive Grösse $\Theta$ des Körpers ausgedrückt werden.
\end{enumerate}
Denn sind $\omega,\omega_1,\ldots,\omega_{n-1}$ die zu $\omega$ conjugirten Zahlen, ebenso $\Theta,\Theta_1,\ldots,\Theta_{n-1}$ die zu $\Theta$ conjugirten, so ist
\[
        \Phi(t)=(t-\Theta)(t-\Theta_1)\cdots(t-\Theta_{n-1}),
\]
und
\[
        \Phi(t)\left(
        \frac{\omega}{t-\Theta}+
        \frac{\omega_1}{t-\Theta_1}+\cdots+
        \frac{\omega_{n-1}}{t-\Theta_{n-1}}
        \right)=\Psi(t)
\]
eine ganze rationale Function $n-1$ten Grades von $t$ in $\Omega$. Setzt man $t=\Theta$, so ergiebt sich
\[
        \omega=\frac{\Psi(\Theta)}{\Phi'(\Theta)}.
\]
Es ist hiernach der Körper $\Omega(\Theta)$ mit dem Körper $\Omega(\alpha)$ identisch.

\begin{enumerate}[label=\arabic*.,leftmargin=2em,start=4]
\item Ist $\Theta$ nicht primitiv, so kann nicht jede Grösse in $\Omega(\alpha)$ rational durch $\Theta$ ausgedrückt werden. Der Körper $\Omega(\Theta)$ ist ein Theiler des Körpers $\Omega(\alpha)$ und der Grad von $\Omega(\Theta)$ ist ein Theiler des Grades von $\Omega(\alpha)$.
\end{enumerate}

Der Körper $\Omega(\alpha)$ heisst primitiv, wenn er ausser den Grössen in $\Omega$ keine imprimitiven Grössen enthält, imprimitiv, wenn er noch andere imprimitive Grössen enthält. Ein Körper, dessen Grad eine Primzahl ist, ist nothwendig primitiv.

Hat $\Omega(\alpha)$ einen von $\Omega$ verschiedenen Theiler $\Omega'$, der seinerseits $\Omega$ als Theiler enthält, und ist $\beta$ eine Grösse, die dem Körper $\Omega'$ aber nicht $\Omega$ angehört, so ist $\Omega(\beta)$ ein algebraischer Körper über $\Omega$ und zugleich ein Theiler von $\Omega'$ und von $\Omega(\alpha)$. Fahren wir so fort, so kommen wir endlich zum Abschluss, da die Grade der Körper wachsen und doch kleiner als $n$ bleiben. Daraus folgt:
\begin{enumerate}[label=\arabic*.,leftmargin=2em,start=5]
\item Jeder Theiler von $\Omega(\alpha)$, der den Körper $\Omega$ enthält, ist ein algebraischer Körper $\Omega(\beta)$ über $\Omega$. Der Grad von $\Omega(\beta)$ ist ein Theiler des Grades von $\Omega(\alpha)$, und wenn beide Körper den gleichen Grad haben, so sind sie identisch.
\item Ein algebraischer Körper über $\Omega$ ist primitiv, wenn er ausser $\Omega$ und sich selbst keinen Körper über $\Omega$ zum Theiler hat.
\end{enumerate}

\sect{145}{Normalkörper. Galois'sche Resolvente}

Ist $\Omega(\alpha)$ ein algebraischer Körper, so sind die conjugirten Körper
\[
        \Omega(\alpha),\Omega(\alpha_1),\ldots,\Omega(\alpha_{m-1})
\]
alle von gleichem Grade, und wenn also einer von ihnen im anderen enthalten ist, so sind beide identisch. Ein Körper, der mit allen seinen conjugirten Körpern identisch ist, heisst ein Normalkörper. Die Normalkörper heissen daher auch Galois'sche Körper.

Wenn ein Körper $\mu$ten Grades $\Omega(\rho)$ die Eigenschaft hat, dass die zu $\rho$ conjugirten Zahlen $\rho_1,\rho_2,\ldots,\rho_{\mu-1}$ alle in $\Omega(\rho)$ enthalten sind, so ist es ein Normalkörper.

Wir wollen eine Gleichung eine Normalgleichung nennen, wenn sie irreducibel ist und die Eigenschaft hat, dass alle ihre Wurzeln rational in $\Omega$ durch eine von ihnen ausgedrückt werden können. Dann ist ein primitives Element eines Normalkörpers $\mu$ten Grades Wurzel einer Normalgleichung $\mu$ten Grades. Aber es gilt auch das Umgekehrte, dass, wenn $\rho$ eine Wurzel einer Normalgleichung ist, $\Omega(\rho)$ ein Normalkörper gleichen Grades ist.

Man kann nun auf folgendem Wege beliebige algebraische Körper auf Normalkörper zurückführen. Es sei
\[
\tag{1}
        F(x)=0
\]
eine beliebige reducible oder irreducible Gleichung in $\Omega$ vom $m$ten Grade, die keine mehrfachen Wurzeln habe. Ihre Wurzeln seien
\[
\tag{2}
        \alpha,\alpha_1,\ldots,\alpha_{m-1}.
\]
Man erhält daraus $m$ algebraische Körper
\[
\tag{3}
        \Omega(\alpha),\Omega(\alpha_1),\ldots,\Omega(\alpha_{m-1}).
\]
Wir nennen den aus allen Grössen (2) abgeleiteten Körper über $\Omega$
\[
\tag{4}
        N=\Omega(\alpha,\alpha_1,\ldots,\alpha_{m-1})
\]
den Galois'schen Körper der Gleichung $F(x)=0$. Ist $F(x)$ irreducibel, so sind die Körper (3) die mit $\Omega(\alpha)$ conjugirten Körper. In diesem Falle soll $N$ die Norm eines jeden der Körper $\Omega(\alpha),\Omega(\alpha_1),\ldots$ heissen.

Es bleibt nachzuweisen, dass $N$ ein Normalkörper ist. Wir wählen ein primitives Element $\rho$ des Körpers $N$ und bezeichnen $N$ auch durch $\Omega(\rho)$. Ist $\mu$ der Grad von $N$, so genügt $\rho$ einer irreducibeln Gleichung $\mu$ten Grades
\[
\tag{5}
        g(t)=0.
\]
Die eine Wurzel $\rho$ ist eine rationale Function der $\alpha,\alpha_1,\ldots,\alpha_{m-1}$. Schreiben wir
\[
        \rho=R(\alpha,\alpha_1,\ldots,\alpha_{m-1}),
\]
und bilden alle verschiedenen Anordnungen der Ziffern $0,1,\ldots,m-1$, so erhalten wir Grössen
\[
\tag{7}
        \rho,\rho',\rho'',\ldots
\]
von der Gestalt $R(\alpha_0,\alpha_1,\ldots)$ mit permutirten Argumenten. Durch jede Vertauschung der Wurzeln werden diese Functionen nur unter einander permutirt. Bilden wir also
\[
        G(t)=(t-\rho)(t-\rho')(t-\rho'')\cdots,
\]
so sind die Coefficienten symmetrische Functionen der Wurzeln von $F(x)$, also Functionen von $t$ in $\Omega$. Nun haben $G(t)$ und $g(t)$ eine Wurzel gemein. Da aber $g(t)$ irreducibel ist, so muss $G(t)$ durch $g(t)$ theilbar sein; also sind alle Wurzeln von $g(t)$ in $N$ enthalten. Mithin ist $N$ ein Normalkörper.

Jede Gleichung $g(t)=0$ heisst eine Galois'sche Resolvente der Gleichung $F(x)=0$. Sie ist dadurch definirt, dass $g(t)$ irreducibel ist, dass alle Wurzeln von $F(x)$ rational durch eine Wurzel $\rho$ von $g(t)$ ausdrückbar sind, und dass eine Wurzel von $g(t)$ rational durch die Wurzeln von $F(x)$ ausdrückbar ist.

Jede Galois'sche Resolvente ist eine Normalgleichung, und eine Normalgleichung ist ihre eigene Galois'sche Resolvente.

\sect{146}{Die Substitutionen eines Normalkörpers}

Es sei jetzt $\Omega(\rho)$ ein Normalkörper $\mu$ten Grades und $\rho$ eine seiner primitiven Zahlen, ferner
\[
\tag{1}
        g(t)=0
\]
die irreducible Gleichung $\mu$ten Grades, deren eine Wurzel $t=\rho$ ist. Die zu $\rho$ conjugirten Elemente seien
\[
\tag{2}
        \rho,\rho_1,\rho_2,\ldots,\rho_{\mu-1}.
\]
Da nach der Definition des Normalkörpers die Grössen (2) alle in $\Omega(\rho)$ enthalten sind, so können wir $\Theta_1(t),\Theta_2(t),\ldots,\Theta_{\mu-1}(t)$ als ganze rationale Functionen in $\Omega$ höchstens vom Grade $\mu-1$ so bestimmen, dass
\[
\tag{3}
        \rho_1=\Theta_1(\rho),\quad
        \rho_2=\Theta_2(\rho),\quad\ldots,\quad
        \rho_{\mu-1}=\Theta_{\mu-1}(\rho)
\]
wird. Ist $\omega$ eine beliebige Grösse in $\Omega(\rho)$, so kann man die mit $\omega$ conjugirten Grössen so darstellen:
\[
\tag{4}
        \omega=\varphi(\rho),\quad
        \omega_1=\varphi(\rho_1),\ldots,\quad
        \omega_{\mu-1}=\varphi(\rho_{\mu-1}),
\]
worin $\varphi(t)$ eine rationale Function in $\Omega$ ist.

Da nun $g(t)$ irreducibel ist, so gilt:
\begin{enumerate}[label=\arabic*.,leftmargin=2em]
\item Wenn eine rationale Function $\Phi(t)$ in $\Omega$ eine Wurzel mit $g(t)$ gemeinsam hat, so verschwinden alle conjugirten Grössen
\[
        \Phi(\rho),\Phi(\rho_1),\ldots,\Phi(\rho_{\mu-1}).
\]
\end{enumerate}

Wird in einer der Functionen $\Theta_k(t)$, durch die nach (3) die Wurzeln von (1) ausgedrückt sind, $\rho$ durch eine andere Wurzel $\rho_h$ ersetzt, so entsteht daraus wieder eine der Wurzeln. Die Reihe der Grössen
\[
\tag{5}
        \rho_h,\Theta_1(\rho_h),\Theta_2(\rho_h),\ldots,\Theta_{\mu-1}(\rho_h)
\]
stimmt, von der Anordnung abgesehen, mit der Reihe
\[
\tag{6}
        \rho,\Theta_1(\rho),\Theta_2(\rho),\ldots,\Theta_{\mu-1}(\rho)
\]
überein. Denn wenn zwei der Grössen (5) gleich wären, so folgte aus
\[
\tag{7}
        \Theta_i(\rho_h)=\Theta_k(\rho_h)
\]
nach Satz 1
\[
\tag{8}
        \Theta_i(\rho)=\Theta_k(\rho),
\]
was für $i\ne k$ nicht möglich ist.

\begin{enumerate}[label=\arabic*.,leftmargin=2em,start=2]
\item Vertauscht man $\rho$ mit $\rho_h$, so geht zugleich jede mit $\rho$ conjugirte Grösse in eine bestimmte andere über, und niemals zwei verschiedene in dieselbe.
\end{enumerate}

Wenn wir in allen Functionen von $\rho$ statt $\rho$ eine andere Wurzel $\rho_a$ setzen, so führen wir eine Substitution aus. Wir bezeichnen diese Substitution mit
\[
        \delta_a=(\rho,\rho_a),\qquad a=0,1,2,\ldots,\mu-1,
\]
wobei unter $\delta_0$ oder $\delta$ die identische Substitution $(\rho,\rho)$ verstanden wird.

Wenn wir in einer beliebigen der Wurzeln
\[
        \rho_h=\Theta_h(\rho)
\]
die Substitution $\delta_a$ ausführen, so geht $\rho_h$ in eine andere Wurzel $\rho_k$ über, die bestimmt ist durch
\[
        \rho_k=\Theta_h(\rho_a).
\]
Eine beliebige Grösse $\omega=\varphi(\rho)$ des Körpers $\Omega(\rho)$ geht durch die Substitution $\delta_a$ in $\omega_a=\varphi(\rho_a)$ über. Drücken wir $\omega$ durch $\rho_h$ und $\omega_a$ durch $\rho_k$ aus, so sieht man, dass $\omega$ in $\omega_a$ übergeht durch die Substitution
\[
\tag{9}
        \delta_a=(\rho_h,\rho_k)=(\rho_h,\Theta_h\Theta_a(\rho)).
\]
Hierin ist $\rho_h$ eine beliebige Wurzel von $g(t)$ und das zugehörige $\rho_k$ ist durch $\delta_a$ bestimmt.

\begin{enumerate}[label=\arabic*.,leftmargin=2em,start=3]
\item Jede Substitution $\delta_a$ kann in der Form $(\rho_h,\rho_k)$ dargestellt werden, worin entweder $\rho_h$ oder $\rho_k$ eine beliebig gegebene der $\mu$ Wurzeln $\rho$ ist, während die andere durch $\delta_a$ völlig bestimmt ist.
\end{enumerate}

Die Anzahl aller von einander verschiedenen Substitutionen $\delta_a$ ist, die identische Substitution mitgerechnet, gleich dem Grade des Körpers $\Omega(\rho)$. Jede dieser Substitutionen führt die Gesammtheit der Grössen des Körpers $\Omega(\rho)$ in sich selbst über, so dass jede in eine bestimmte andere übergeht, und niemals zwei verschiedene Grössen in die gleiche.

Wir nennen daher die $\delta_a$ die Substitutionen des Körpers $\Omega(\rho)$. Wenn $\omega=\varphi(\rho)$ im Körper $\Omega(\rho)$ ungeändert bleibt, wenn $\rho$ durch $\rho_a$ ersetzt wird, wenn also
\[
        \varphi(\rho)=\varphi(\rho_a)
\]
ist, so sagen wir, $\omega$ erlaubt oder gestattet die Substitution $(\rho,\rho_a)$ oder $\delta_a$.

Die Grössen, die ausser der identischen Substitution keine Substitution $\delta_a$ gestatten, sind die primitiven Elemente des Körpers $\Omega(\rho)$.
\begin{enumerate}[label=\arabic*.,leftmargin=2em,start=4]
\item Eine Grösse, die alle Substitutionen $\delta_a$ gestattet, ist nothwendig ein Element von $\Omega$.
\end{enumerate}
Denn eine solche Grösse ist mit allen ihren conjugirten Grössen identisch, und genügt also nach §.\,144 einer Gleichung ersten Grades in $\Omega$.

\sect{147}{Zusammensetzung der Substitutionen}

Wenn wir in irgend einer Function von $\rho$ die Wurzel $\rho$ zuerst durch $\rho_a$ und dann $\rho_a$ durch $\rho_b$ ersetzen, so ist der Erfolg derselbe, als ob wir gleich von vornherein $\rho$ mit $\rho_b$ vertauscht hätten. Setzen wir also
\[
\tag{1}
        \delta=(\rho,\rho_a),\qquad
        \delta'=(\rho_a,\rho_b),\qquad
        \delta''=(\rho,\rho_b),
\]
so ist es für das Ergebniss gleichgültig, ob wir zuerst die Substitution $\delta$ und dann die Substitution $\delta'$ ausführen, oder ob wir für einmal die Substitution $\delta''$ machen.

\begin{enumerate}[label=\arabic*.,leftmargin=2em]
\item Wir nennen daher $\delta''$ aus $\delta$ und $\delta'$ componirt oder zusammengesetzt, und bezeichnen diese Beziehung durch die symbolische Gleichung
\[
\tag{2}
        \delta\delta'=\delta''.
\]
\end{enumerate}

Nach Satz 3 des vorigen Paragraphen lassen sich $\rho_a,\rho_b$ so auswählen, dass $\delta,\delta'$ zwei beliebig gegebene unter den $\mu$ Substitutionen sind; $\delta''$ ist dann völlig bestimmt. Ebenso ist aber auch, wenn $\delta$ und $\delta''$ gegeben sind, $\delta'$ eindeutig bestimmt, und wenn $\delta'$ und $\delta''$ gegeben sind, ist $\delta$ eindeutig bestimmt.

\begin{enumerate}[label=\arabic*.,leftmargin=2em,start=2]
\item Von den durch die symbolische Gleichung $\delta\delta'=\delta''$ verbundenen drei Substitutionen des Körpers $\Omega(\rho)$ können irgend zwei beliebig gegeben sein, während die dritte dadurch eindeutig bestimmt ist.
\end{enumerate}

Wir haben die Composition der Substitutionen symbolisch durch das Zeichen der Multiplication ausgedrückt. Es ist hier aber wohl auf die Reihenfolge der Factoren oder Componenten zu achten, weil $\delta\delta'$ von $\delta'\delta$ verschieden sein kann; bei der die Composition ausdrückenden symbolischen Multiplication gilt also nicht das commutative Gesetz der gewöhnlichen Multiplication. Wohl aber gilt das associative Gesetz:

\begin{enumerate}[label=\arabic*.,leftmargin=2em,start=3]
\item Sind $\delta,\delta',\delta''$ irgend drei der $\mu$ Substitutionen des Körpers $\Omega(\rho)$, so ist
\[
\tag{3}
        (\delta\delta')\delta''=\delta(\delta'\delta'').
\]
\end{enumerate}

Der Beweis ist sehr einfach; denn nach Satz 3, §.\,146 können wir $\rho_a,\rho_b,\rho_c$ so bestimmen, dass
\[
\tag{4}
        \delta=(\rho,\rho_a),\qquad
        \delta'=(\rho_a,\rho_b),\qquad
        \delta''=(\rho_b,\rho_c),
\]
und dann ist
\[
\delta\delta'=(\rho,\rho_b),\qquad
(\delta\delta')\delta''=(\rho,\rho_c),
\]
\[
\delta'\delta''=(\rho_a,\rho_c),\qquad
\delta(\delta'\delta'')=(\rho,\rho_c).
\]
Wir bezeichnen daher kurz die aus den drei Substitutionen zusammengesetzte Substitution mit $\delta\delta'\delta''$, und können überhaupt aus beliebig vielen Componenten eine bestimmte Substitution zusammensetzen, indem wir nach einander immer zwei benachbarte unter den Componenten zu einer einzigen Substitution vereinigen, bis das ganze symbolische Product sich auf eine einzige Substitution zusammengezogen hat. Denn wir können setzen
\[
        \delta=(\rho,\rho_a),\quad
        \delta'=(\rho_a,\rho_b),\quad
        \delta''=(\rho_b,\rho_c),\quad
        \delta'''=(\rho_c,\rho_e),
\]
und erhalten immer
\[
        (\rho,\rho_a)(\rho_a,\rho_b)(\rho_b,\rho_c)(\rho_c,\rho_e)
        =(\rho,\rho_e),
\]
was offenbar auf eine beliebige Anzahl von Componenten ausgedehnt werden kann.


\end{document}
